\section{Results}

\subsection{Token-ID transcription}

Table~\ref{tab:development} shows the development trajectory.  Direct token-ID
prediction is viable but not solved: the selected E1 model improves strict sequence
accuracy by roughly an order of magnitude over the initial autoregressive decoder,
yet most complete sequences still contain at least one error.  The positional
intervention changes pairwise label order from near chance to approximately 0.99,
supporting the claim that the emitted sequence is visually ordered.

\begin{table}[t]
\centering
\caption{Signer-held-out development results for seed 23.  E1 is the selected
token-ID \system{}; folds 4--6 were used for development and are not prospective.}
\label{tab:development}
\begin{tabular}{rrrrr}
\toprule
Fold & Initial AR exact (\%) & E1 exact (\%) & E1 edit (\%) & Pair order (\%) \\
\midrule
4 & 2.2 & 18.1 & 63.0 & 98.6 \\
5 & 2.3 & 19.8 & 61.9 & 98.9 \\
6 & 1.8 & 16.9 & 58.5 & 98.6 \\
\bottomrule
\end{tabular}
\end{table}

\begin{table}[t]
\centering
\caption{Prospectively registered signer-held-out confirmation.  The CIF column
is the frozen historical comparator, not a candidate \system{} architecture.}
\label{tab:confirmation}
\begin{tabular}{rrrrrrrr}
\toprule
Fold & CIF exact & E1 exact & E1 edit & Order & $\Delta_{\mathrm{perm}}$ & W/L/T & $p$ \\
\midrule
7 & \FoldSevenCIFExact & \FoldSevenExact & \FoldSevenEdit & \FoldSevenOrder &
    \FoldSevenPermutationDrop & \FoldSevenWLT & \FoldSevenPValue \\
8 & \FoldEightCIFExact & \FoldEightExact & \FoldEightEdit & \FoldEightOrder &
    \FoldEightPermutationDrop & \FoldEightWLT & \FoldEightPValue \\
\bottomrule
\end{tabular}
\end{table}

Under the registered conjunctive decision rule, H1 \HOneOutcome, H2
\HTwoOutcome, and H3 \HThreeOutcome.  Overall confirmation therefore
\OverallOutcome.  These statements are populated automatically from the frozen
evaluation artifacts and are not manually reinterpreted after either fold.

\subsection{Can \Gemma{} repair token errors?}

Table~\ref{tab:gemma-tolerance} answers this question for the current isolated-label
regime.  At zero corruption, \Gemma{} preserves exact content in 95\% of rows and
keeps chrF near 100.  At every tested nonzero rate, exact accuracy is already zero
before correction and remains zero afterward.  chrF decreases rather than improves.
Thus, no nonzero error tolerance is demonstrated: there is no tested nonzero
corruption level at which the language stage recovers exact content.  Because the
short sequences quantize every nonzero nominal rate to at least one substitution,
this experiment does not locate a continuous threshold between 0 and 0.10.

\begin{table}[t]
\centering
\caption{\Gemma{} 3n E2B correction audit on 120 isolated-label rows per
condition.  chrF is on a 0--100 scale.  On these short sequences, every nonzero
nominal rate produces at least one substitution.}
\label{tab:gemma-tolerance}
\begin{tabular}{rrrrrr}
\toprule
Substitution rate & Exact before & Exact after & chrF before & chrF after & $\Delta$chrF \\
\midrule
0.00 & 1.000 & 0.950 & 100.00 & 99.16 & -0.84 \\
0.10 & 0.000 & 0.000 & 36.67 & 34.58 & -2.10 \\
0.20 & 0.000 & 0.000 & 36.67 & 34.58 & -2.10 \\
0.30 & 0.000 & 0.000 & 36.67 & 34.58 & -2.10 \\
0.40 & 0.000 & 0.000 & 34.29 & 32.35 & -1.94 \\
0.50 & 0.000 & 0.000 & 23.30 & 21.85 & -1.45 \\
\bottomrule
\end{tabular}
\end{table}

This is not evidence that language models cannot correct sign transcription in
general.  A substitution such as ``photo'' $\rightarrow$ ``green'' in an isolated
label is compatible with no surrounding syntax or semantics that identifies the
original word.  The appropriate next experiment is therefore not another prompt
sweep on LSA64, but a pre-registered sentence corpus in which controlled token
errors leave genuine contextual redundancy.
